\clearpage
\begingroup
\fancyhead[L]{\small REMARKS}
\renewcommand{\headrulewidth}{0pt}
\SourceLeaf{607}
\needspace{6\baselineskip}
\pdfbookmark[1]{Remarks on the perturbations of Pallas}{pallas-bemerkungen}
\begin{center}\large REMARKS.\end{center}
\begingroup
\newcommand{\BIntroJ}{\text{\mdseries\AstroSymbols ♃}}
\newcommand{\BIntroP}{\text{\mdseries\AstroSymbols ⚴}}

\par\noindent\hspace*{\parindent}What survives in the posthumous papers from \textsc{Gauss}'s investigations into the perturbations of Pallas consists almost entirely of
numerical calculations without explanatory text. They are written on a great many loose slips,
which fill five archival boxes. The calculation of the oppositions and the comparison with the
observations were, however, often carried out by \textsc{Gauss} in his notebooks as well. The sheer amount of
material made it much harder to sort. It was finally possible, particularly by studying
the correspondence at the same time, to put the whole in coherent order.

\par\noindent\hspace*{\parindent}The passages from the correspondence (pp.~413—438) placed first above
make clear the historical course of \textsc{Gauss}'s investigations. Their connection with the text
printed later calls for the following remarks:

\medskip
\par\noindent\hspace*{\parindent}While calculating the perturbations of Ceres, \textsc{Gauss} conceived the plan of preparing a table
to facilitate expanding the disturbing function for all the minor planets discovered by then.
He wrote about it to \textsc{Olbers} on 3~September, 7~and 29~October 1805. There is also
a set of entries in Notebook B\,e on the method \textsc{Gauss} intended to use.
The table, whose calculation \textsc{Bessel} helped with, was completed and is among the posthumous papers.
\textsc{Gauss} then used other methods when calculating the perturbations of Pallas,
so the table was never used. It was therefore omitted from this edition.

\par\noindent\hspace*{\parindent}The letters to \textsc{Olbers} from January 1806 to October 1810 concern preliminary work on calculating
the perturbations. They deal partly with how precisely the observations can be represented
using the method of least squares, and partly with the analytical
foundations for expanding the disturbing function. At the end of October 1810, \textsc{Gauss} then began the first calculation
of the special perturbations of Pallas (printed on p.~473), which is the subject of the letters from
November 1810. The letter to \textsc{Olbers} of 13~December 1810 concerns the second calculation of the
special perturbations (p.~483). Those of August and October 1811 concern the first calculation of the general
perturbations (p.~489). In April 1812, \textsc{Gauss} began the second calculation of the general perturbations,
which was completed in July 1813. For the period from March 1812 to January 1813, the letter to
\textsc{Bessel} of 5~May 1812 and the letters from \textsc{Olbers} to \textsc{Gauss} are the only ones that could be printed, since the corresponding
letters from \textsc{Gauss} to \textsc{Olbers} have been lost. Among them is the one
in which \textsc{Gauss}, as in the letter to \textsc{Bessel} just mentioned, reports his discovery of the rational ratio
of the periods of revolution of Pallas and Jupiter. \textsc{Gauss} thus first made the discovery before
carrying out the second calculation of the general perturbations, in March or early April 1812,
when he had yet to calculate with the small divisor \(18\BIntroJ-7\BIntroP\). It remains unclear (but see
the remarks below on p.~604) how he found that rational ratio in this particular case. The notes
\par\smallskip\noindent VII.\hfill 76\par
\noindent\SourceLeaf{608}on the subject that were found in the posthumous papers and are printed on p.~558 date from later, in any case
\textls[120]{after} at least a large part of the second calculation had been carried out.

\par\noindent\hspace*{\parindent}\textsc{Gauss} took special care in comparing his results with the observations,
and, in particular, in determining the normal elements which, in conjunction with the calculated perturbations, represented the
observations as closely as possible. This gives a very clear picture of
how closely the observations agree with the theory. Most of the differences between the two
had already been published by \textsc{Gauss} himself (see vol.~VI). They are also found in the comparisons
with the oppositions and in the passages from the letters printed above. The comparisons with the oppositions
printed above alongside the determination of the perturbations are accompanied in the posthumous papers by
a great many further comparisons, including some for later oppositions. Some take Saturn's perturbations into account
and some omit them. Some use the old value of Jupiter's mass and others use an improved value.
The most complete comparison with the first 10 or 12 oppositions (1802 to 1817) survives in Notebook
Bf. It is omitted here because it offers no new method and its
results can be found in the letter to \textsc{Olbers} printed on p.~430. These comparisons with the oppositions
were made by \textsc{Gauss} both with his special perturbations, which he seems to have continued up to 1816,
and with the general perturbations. In calculating the magnitudes of the perturbations from
the expressions for the general perturbations, \textsc{Westphal}, \textsc{Encke}, and \textsc{Nicolai} gave him substantial
assistance. While the tables (p.~572) were still unfinished, groups of perturbation equations
were collected on small slips according to their magnitude and called »Manipel«. Each Manipel
(apart from the last) contained ten perturbation terms. For~example, the first Manipel for the
longitude of perihelion are as follows:

\needspace{19\baselineskip}
\begin{center}\textit{Longitude of perihelion.}\end{center}
\begingroup\small
\setlength{\arraycolsep}{3pt}
\begin{center}\textit{Manipel 1.}\end{center}
\[
\begin{array}{ccc}
\left[\begin{gathered}\text{\itshape Perturbation}\\\text{\itshape equation no.}\end{gathered}\right]&
\text{\itshape [Argument]}&
\left[\begin{gathered}\text{\itshape Log. of}\\\text{\itshape coefficient}\end{gathered}\right]\\
\mathit{27}&\mathit{2}\BIntroJ-\BIntroP+\mathit{236^0\,43'\,49''}&\mathit{2{,}88715}\\
\mathit{66}&\mathit{5}\BIntroJ-\mathit{2}\BIntroP+\mathit{303\ \,26\ \,18}&\mathit{2{,}81492}\\
\mathit{40}&\mathit{3}\BIntroJ-\BIntroP+\mathit{94\ \,47\ \,46}&\mathit{2{,}72535}\\
\mathit{52}&\mathit{4}\BIntroJ-\mathit{2}\BIntroP+\mathit{288\ \,15\ \,4}&\mathit{2{,}16040}\\
\mathit{39}&\mathit{3}\BIntroJ-\mathit{2}\BIntroP+\mathit{260\ \,27\ \,22}&\mathit{1{,}95730}\\
\mathit{26}&\mathit{2}\BIntroJ-\mathit{2}\BIntroP+\mathit{342\ \,9\ \,24}&\mathit{1{,}94234}\\
\mathit{28}&\mathit{2}\BIntroJ\hspace{2em}+\mathit{101\ \,42\ \,47}&\mathit{1{,}91320}\\
\mathit{100}&\mathit{8}\BIntroJ-\mathit{3}\BIntroP+\mathit{168\ \,38\ \,33}&\mathit{1{,}90406}\\
\mathit{53}&\mathit{4}\BIntroJ-\BIntroP+\mathit{106\ \,12\ \,53}&\mathit{1{,}88715}\\
\mathit{14}&\BIntroJ\hspace{2em}+\mathit{291\ \,57\ \,59}&\mathit{1{,}85959}
\end{array}
\]
\begin{center}\textit{Manipel 2.}\end{center}
\[
\begin{array}{ccc}
\left[\begin{gathered}\text{\itshape Perturbation}\\\text{\itshape equation no.}\end{gathered}\right]&
\text{\itshape [Argument]}&
\left[\begin{gathered}\text{\itshape Log. of}\\\text{\itshape coefficient}\end{gathered}\right]\\
\mathit{78}&\mathit{6}\BIntroJ-\mathit{2}\BIntroP+\mathit{137^0\,59'\,40''}&\mathit{1{,}77929}\\
\mathit{13}&\BIntroJ-\BIntroP+\mathit{309\ \,34\ \,12}&\mathit{1{,}76702}\\
\mathit{1}&\hspace{1.5em}\BIntroP+\mathit{321\ \,28\ \,57}&\mathit{1{,}75857}\\
\mathit{89}&\mathit{7}\BIntroJ-\mathit{3}\BIntroP+\mathit{334\ \,36\ \,27}&\mathit{1{,}75507}\\
\mathit{38}&\mathit{3}\BIntroJ-\mathit{3}\BIntroP+\mathit{340\ \,39\ \,45}&\mathit{1{,}68732}\\
\mathit{77}&\mathit{6}\BIntroJ-\mathit{3}\BIntroP+\mathit{320\ \,40\ \,29}&\mathit{1{,}59337}\\
\mathit{65}&\mathit{5}\BIntroJ-\mathit{3}\BIntroP+\mathit{306\ \,34\ \,54}&\mathit{1{,}57175}\\
\mathit{122}&\mathit{10}\BIntroJ-\mathit{4}\BIntroP+\mathit{18\ \,48\ \,20}&\mathit{1{,}54904}\\
\mathit{41}&\mathit{3}\BIntroJ\hspace{2em}+\mathit{87\ \,0\ \,49}&\mathit{1{,}50939}\\
\mathit{12}&\BIntroJ-\mathit{2}\BIntroP+\mathit{242\ \,28\ \,11}.&\mathit{1{,}43467}
\end{array}
\]
\endgroup
\begin{center}\textit{etc.}\end{center}

\par\noindent For each Manipel the individual perturbation terms were calculated and then the amounts of the
Manipel were added to find the total perturbations of each element. Special perturbations for later
years were also calculated by \textsc{Encke} (see~Berliner Astronomisches Jahrbuch für 1837—38).

\par\noindent\hspace*{\parindent}The first investigation of Jupiter's mass is printed on p.~561. The course of \textsc{Gauss}'s work
on this is shown by the correspondence on pp.~424—429. It is striking how accurately \textsc{Gauss} ultimately gives the value
of this quantity in his letter to \textsc{Olbers} of 24~July 1816 as \(\frac{1}{1030}\).

\par\noindent\hspace*{\parindent}The letters concerning the tables of perturbations by Jupiter printed on p.~572 date from
\SourceLeaf{609}February 1817 to March 1818. The first plan for such a table is already mentioned in the letter to \textsc{Olbers} of
8~April 1813, in which \textsc{Gauss} also expresses his intention to have the perturbations by Saturn and
Mars calculated.

\par\noindent\hspace*{\parindent}Hints about the development of his method, which we must take to have begun with the »Exposition d’une nouvelle méthode
etc.«, are found in the letters of 12~December 1813 and 25~September
1814. That he began the work intending to compete for the prize offered by the French Academy
is already clear from his writing it in French. How the
astronomical community's hopes that the prize would go to \textsc{Gauss} persisted for many years
and were ultimately disappointed is plain enough from the correspondence.
\endgroup
\endgroup
